\section{Вступ і мотивація}

Стрімінгові музичні платформи генерують значні обсяги даних про музичні композиції, поведінку користувачів та характеристики аудіоконтенту. Такі дані є цінним джерелом інформації для задач інтелектуального аналізу даних, оскільки дозволяють досліджувати закономірності популярності музики, взаємозв'язки між жанрами та аудіофічами, а також виконувати класифікацію й кластеризацію треків.

Платформа Spotify надає набір аудіофіч, які описують композиції з точки зору ритму, енергійності, танцювальності, акустичності, темпу та інших параметрів. Аналіз таких характеристик дозволяє досліджувати структуру музичних даних та визначати фактори, що найбільше впливають на популярність композицій.

Метою цієї роботи є проведення комплексного аналізу набору даних Spotify Tracks Dataset із використанням методів інтелектуального аналізу даних та машинного навчання. У межах дослідження виконано попередню обробку даних, описовий аналіз, кореляційний аналіз, кластеризацію треків та побудову моделей класифікації популярності композицій.

Основними завданнями роботи є:
\begin{enumerate}
  \item Аналіз структури та характеристик набору даних.
  \item Дослідження взаємозв'язків між аудіофічами.
  \item Виявлення груп подібних треків за допомогою кластеризації.
  \item Побудова моделей машинного навчання для прогнозування популярності треків.
  \item Інтерпретація отриманих результатів.
\end{enumerate}

Набір даних містить інформацію про понад 100 тисяч музичних композицій різних жанрів. Це забезпечує достатню репрезентативність вибірки для проведення аналізу та отримання практично значущих результатів.
